\introduction % Структурный элемент: ВВЕДЕНИЕ

\textbf{Актуальность исследования}

Актуальность выпускной квалификационной работы определяется тем, что
прикладные системы, использующие данные сети Bitcoin, нуждаются не только в
доступе к исходным блокам и транзакциям, но и в устойчивом слое
структурированных представлений. Узел Bitcoin Core является каноническим
источником информации о состоянии сети и предоставляет набор RPC-методов для
получения блоков, транзакций, сведений о цепи и UTXO-наборе
\cite{btc_core_rpc, bitcoin_core_rpc_29}. Однако прямой RPC-доступ ориентирован
прежде всего на работу узла как участника сети и не предназначен для
регулярного обслуживания прикладных выборок, построения истории балансов,
фильтрации транзакций по адресам и предоставления внешним клиентам удобного
программного интерфейса.

В практических сценариях это приводит к нескольким проблемам. Во-первых,
внешнему приложению приходится повторно обходить блоки, транзакции, входы и
выходы, чтобы получить производные данные, например текущий набор UTXO адреса
или историю изменения его баланса. Во-вторых, часть запросов к Bitcoin Core
может быть вычислительно дорогой: операции, связанные со сканированием
UTXO-набора или фильтрацией mempool, требуют дополнительной обработки на
стороне клиента и увеличивают задержку ответа. В-третьих, прикладная система
должна учитывать особенности предметной области Bitcoin: наличие канонической
цепи, возможные reorg, временное состояние mempool и необходимость отделять
подтвержденные данные от неподтвержденных \cite{nakamoto2008,
narayanan2016bitcoin, antonopoulos2017mastering}.

Для внешних сервисов, аналитических модулей, explorer-backend и внутренних
инструментов мониторинга такое положение означает рост сложности интеграции.
Разработчику необходимо самостоятельно поддерживать локальные агрегаты,
следить за согласованностью данных при изменении цепи, контролировать прогресс
синхронизации и реализовывать повторяющиеся механизмы доступа к данным. При
увеличении числа запросов это создает дополнительную нагрузку на RPC-интерфейс
узла, затрудняет диагностику ошибок и делает поведение прикладного слоя менее
предсказуемым.

Проблема, рассматриваемая в работе, заключается в отсутствии в базовом узле
Bitcoin Core специализированного слоя индексирования, который одновременно
обеспечивает получение данных от доверенного узла, нормализацию блокчейн
данных, транзакционное сохранение в реляционной базе данных, расчет UTXO и
адресных балансов, обработку mempool и reorg, а также прикладную выдачу данных
через API. Существующие решения, такие как explorer-ориентированные индексеры,
аналитические фреймворки и пакетные ETL-инструменты, решают часть этих задач,
но отличаются назначением, моделью хранения, способом эксплуатации и
предоставляемым программным интерфейсом.

В рамках настоящей работы рассматривается разработка программного комплекса
Bitcoin Blockchain Indexer. Система ориентирована на использование в составе
инфраструктуры, где есть доверенный узел Bitcoin Core, реляционная база данных
PostgreSQL и потребители прикладных данных. Разрабатываемый комплекс получает
информацию от узла, сохраняет блоки, транзакции, входы, выходы, UTXO, текущие
и исторические балансы, поддерживает управление заданиями индексирования,
публикует REST API, предоставляет диагностическую информацию и метрики
наблюдаемости \cite{postgres_docs, prometheus_docs}.

Проверка разработанного решения выполняется через сценарии, отражающие основные
бизнес-функции системы: управление заданиями индексирования, получение данных
об адресных балансах, UTXO, транзакциях, блоках и mempool, обработку
расхождения канонической цепи и восстановление производных состояний после
reorg. Дополнительно в работе рассматриваются замеры производительности,
сравнивающие прямые обращения к Bitcoin Core RPC и чтение из заранее
подготовленного PostgreSQL-индекса. Такой подход позволяет оценить не только
корректность отдельных модулей, но и практический эффект от выделения
специализированного слоя индексирования.

\textbf{Постановка задачи}

Целью выпускной квалификационной работы является проектирование и реализация
системы индексации блокчейн данных, обеспечивающей согласованное хранение и
прикладную выдачу данных сети Bitcoin.

Для достижения поставленной цели необходимо решить следующие задачи:
\begin{enumerate}
    \item проанализировать предметную область индексирования блокчейн данных
    сети Bitcoin, включая модель блоков и транзакций, UTXO-набор, адресные
    балансы, mempool, reorg и существующие подходы к получению прикладных
    представлений;
    \item сформировать требования к системе Bitcoin Blockchain Indexer,
    разделив функциональные требования к индексированию, API и управлению
    заданиями, а также нефункциональные требования к производительности,
    надежности, наблюдаемости и контейнерному запуску;
    \item обосновать архитектурный подход к построению системы, определить
    состав модулей, границы ответственности backend-компонента, роль
    доверенного узла Bitcoin Core и способ хранения критического состояния во
    внешней базе данных;
    \item спроектировать модель данных и схему хранения для блоков,
    транзакций, входов, выходов, UTXO, текущих и исторических адресных
    балансов, заданий индексирования, сведений об узлах и диагностических
    показателей;
    \item реализовать ключевые модули программного комплекса: получение данных
    через RPC, пайплайн индексации, синхронизацию mempool, обработку reorg,
    сервис управления заданиями, прикладной REST API, публикацию метрик и
    конфигурацию контейнерного запуска;
    \item провести проверку качества разработанного решения, включая
    тестирование бизнес-сценариев, проверку корректности UTXO и балансов,
    обработку mempool и reorg, а также сравнительные испытания
    производительности относительно прямого доступа через Bitcoin Core RPC.
\end{enumerate}

\textbf{Объект и предмет исследования}

Объектом исследования являются процессы получения, структурированного хранения
и прикладного использования блокчейн данных сети Bitcoin во внешних
информационных системах. К таким процессам относятся синхронизация данных с
доверенным узлом, подготовка производных представлений, обслуживание запросов
клиентов и поддержание согласованности состояния при изменениях канонической
цепи.

Предметом исследования являются методы и программные средства индексирования
блоков, транзакций, UTXO, адресных балансов и mempool-данных, а также подходы к
построению API и модели хранения, позволяющие снизить зависимость прикладных
клиентов от низкоуровневых RPC-вызовов Bitcoin Core.

\textbf{Методы исследования}

Для выполнения работы использовались методы системного анализа предметной
области, сравнительного анализа существующих программных решений,
формализации требований, проектирования программной архитектуры и модели
данных. При реализации прототипа применялись методы модульного проектирования,
транзакционной обработки данных в реляционной базе, проектирования REST API,
контейнеризации и организации наблюдаемости через метрики.

Для оценки результата использовались сценарное тестирование и интеграционные
проверки основных бизнес-функций: жизненного цикла заданий индексирования,
выдачи данных через API, записи блоков и транзакций, обновления UTXO и
адресных балансов, обработки mempool и отката состояния при reorg. Для оценки
эксплуатационных характеристик применялись сравнительные замеры задержек,
пропускной способности и успешности ответов при обращении к Bitcoin Core RPC и
к API разработанного индексера.

\textbf{Практическая значимость и новизна}

Практическая значимость работы состоит в том, что разработанный программный
комплекс предоставляет прикладной слой доступа к данным Bitcoin поверх
доверенного узла Bitcoin Core. Внешний потребитель может обращаться к
подготовленным представлениям через REST API, получать текущие и исторические
балансы адресов, наборы UTXO, сведения о блоках, транзакциях и mempool, а также
контролировать выполнение заданий индексирования без самостоятельного обхода
сырого RPC-интерфейса.

Новизна работы состоит не в создании альтернативного узла Bitcoin и не в
изменении правил консенсуса, а в объединении в одном программном комплексе
нескольких прикладных механизмов: последовательной индексации канонической
цепи, транзакционного хранения нормализованных данных, расчета текущих и
исторических балансов, разделения подтвержденных и mempool-данных, обработки
reorg, управления заданиями индексирования и публикации эксплуатационных
метрик. Такое сочетание делает систему пригодной для дальнейшего развития в
направлении explorer-backend, аналитических сервисов и внутренних
инструментов мониторинга.

\textbf{Основные результаты работы}

В результате выполнения выпускной квалификационной работы были получены
следующие результаты:
\begin{enumerate}
    \item выполнен анализ предметной области индексирования блокчейн данных
    Bitcoin, особенностей UTXO-модели, mempool, reorg и существующих подходов к
    получению прикладных представлений;
    \item сформулированы функциональные и нефункциональные требования к системе
    Bitcoin Blockchain Indexer, включая требования к API, управлению заданиями,
    производительности, надежности и наблюдаемости;
    \item обоснована архитектура stateless модульного монолита с вынесением
    критического состояния в PostgreSQL и использованием доверенного узла
    Bitcoin Core как источника исходных данных;
    \item спроектирована модель хранения блоков, транзакций, входов, выходов,
    UTXO, адресных балансов, истории балансов, заданий индексирования и
    сведений об узлах;
    \item реализованы ключевые модули индексатора, включая RPC-клиент,
    пайплайн записи блоков, сервис mempool, обработку reorg, jobs API, data
    API, метрики, конфигурацию и контейнерный запуск;
    \item проведены проверки бизнес-функций и сравнительные испытания
    производительности, а исходный код проекта размещен в Git-репозитории
    \cite{project_github}.
\end{enumerate}

\textbf{Ограничения исследования}

Результаты работы получены на прототипе системы и на ограниченном наборе
сценариев, достаточном для подтверждения основных проектных решений. В работе
не рассматриваются вопросы промышленного горизонтального масштабирования,
межкластерной отказоустойчивости, эксплуатации под длительной параллельной
нагрузкой и полного security-аудита. Замеры производительности интерпретируются
в рамках зафиксированного локального стенда, поскольку абсолютные значения
задержек и пропускной способности зависят от аппаратной конфигурации, версии
Bitcoin Core, настроек PostgreSQL и профиля нагрузки.

Указанные ограничения не снижают практической значимости разработанного
решения, однако определяют направления дальнейшего развития: расширение
нагрузочных испытаний, проведение полного end-to-end тестирования с реальным
Bitcoin Core/regtest, усиление эксплуатационного контура безопасности,
оптимизацию запросов и развитие административных инструментов сопровождения.
